\section{Loading local FASTA, FASTQ, and gzip files}
\label{sec:appendix-local-ingestion}
\label{sec:app-local-ingestion}

Use \texttt{bioseq.load} to put a file into an existing table created with
\texttt{USING bioseq}. The file must be visible to the PostgreSQL server.
The destination sequence column determines the alphabet and storage width;
the loader can also put parsed header fields into ordinary typed columns.
For a first exercise without files or network access, use the generators in
Section~\ref{sec:appendix-sequence-generation}.

\subsection{Recognize supported file names}

The base suffix selects the record grammar. The accepted names are:
\begin{itemize}
  \item FASTA: \texttt{.fasta}, \texttt{.fa}, or \texttt{.fna}.
  \item FASTQ: \texttt{.fastq}, \texttt{.fq}, or \texttt{.fnq}.
  \item Either group may end in \texttt{.gz}, \texttt{.bgz}, or
        \texttt{.bgzf}. Suffix matching ignores ASCII letter case.
\end{itemize}
The reader checks gzip magic bytes to identify compression. It does not infer
FASTA versus FASTQ from each record. BGZF is accepted as a gzip stream: the
loader reads from the beginning in order and does not use a BGZF index or
perform random range reads. Local ZIP and tar archives are unsupported.
The NCBI assembly ZIP route belongs to the remote interface described below.

\subsection{Design the destination and map headers}

The destination must have exactly one \texttt{bioseq.seq} column with a bound
type modifier, such as \texttt{seq('dna5',64)}. A sequence column without an
alphabet and width, or a table with several sequence columns, is rejected.
Identity and generated columns are unsupported by this loader. An unmapped
ordinary column receives its default, or \texttt{NULL} if it has no default;
an unmapped \texttt{NOT NULL} column therefore needs a usable default.

Use an alphabet that retains the source symbols. For example, DNA5 retains
\texttt{N}. Loading N-bearing input into DNA4 can shorten the stored sequence
and move later coordinates. The file encoder has permissive, path-dependent
handling of unsupported bytes; it is not a strict alphabet validator.
Typed SQL casts validate their input separately. Validate files before
loading when preserving symbols and coordinates is required.

The public call has four arguments and returns the number of complete
sequence records inserted:
\par\noindent\begin{minipage}{\linewidth}
\begin{lstlisting}[style=bioseq-appendix-sql]
bioseq.load(
    target_table regclass,
    file_path text,
    header_format text DEFAULT 'auto',
    columns bioseq.column_mapping[] DEFAULT NULL
) RETURNS bigint
\end{lstlisting}
\end{minipage}

Each mapping pairs a destination column name with a parsed header field.
For a file using UniProt-style headers, the following example stores the
accession and organism and supplies a date from the column default. The
example path refers to a file installed on the database server.
\par\noindent\begin{minipage}{\linewidth}
\begin{lstlisting}[style=bioseq-appendix-sql]
CREATE TABLE imported_reads (
    accession text,
    organism_name text,
    collected_on date DEFAULT CURRENT_DATE,
    seq bioseq.seq('dna5', 64) NOT NULL
) USING bioseq;

SELECT bioseq.load(
    'imported_reads'::regclass,
    '/srv/import/reads.fasta.gz',
    'uni_prot',
    ARRAY[
        ('accession', 'accession')::bioseq.column_mapping,
        ('organism_name', 'organism_name')::bioseq.column_mapping
    ]
);

SELECT accession, organism_name, collected_on,
       bioseq.len(seq) AS bases
FROM imported_reads
ORDER BY accession;
\end{lstlisting}
\end{minipage}

Header syntax and sequence alphabet are separate choices. The available
header formats are \texttt{auto}, \texttt{plain}, \texttt{uni\_prot},
\texttt{uni\_ref}, \texttt{uni\_parc}, \texttt{uni\_mes}, \texttt{spd},
\texttt{ncbi}, \texttt{ncbi\_gi}, \texttt{ncbi\_pdb}, \texttt{refseq},
\texttt{genbank}, \texttt{gisaid}, and \texttt{imgt}.
With \texttt{auto}, a bounded sample of decompressed headers selects one
format for the file. This does not classify every header independently or
validate every dialect convention. For a reproducible import contract,
specify the format and mappings explicitly; provider conventions can change.

There is also a raw-header shortcut. With \texttt{auto}, no explicit mappings,
and exactly one ordinary \texttt{text} column beside the sequence column,
that text column receives the header without the FASTA or FASTQ marker.
The loader then skips dialect detection. For example:
\par\noindent\begin{minipage}{\linewidth}
\begin{lstlisting}[style=bioseq-appendix-sql]
CREATE TABLE incoming_reads (
    header text,
    seq bioseq.seq('dna5', 64) NOT NULL
) USING bioseq;

SELECT bioseq.load(
    'incoming_reads'::regclass,
    '/srv/import/run-042.fastq.gz'
);
\end{lstlisting}
\end{minipage}

Mapped fields pass through PostgreSQL type conversion. Dates, numeric ranges,
enums, and domains can therefore reject unsuitable header text. A failure
reports the record and column context. The loader checks \texttt{NOT NULL}
and table \texttt{CHECK} constraints, but it does not run ordinary
\texttt{INSERT} triggers, maintain ordinary secondary indexes, or perform
their unique, exclusion, and foreign-key checks. Use a staging table and a
subsequent ordinary \texttt{INSERT} to apply supported insertion rules.
\texttt{USING bioseq} rejects foreign keys and \texttt{AFTER ROW} triggers;
those policies need a destination table access method that supports them.

\subsection{Load, build indexes, and search}

For a bulk import, turn off insertion-time BioSeqDB index maintenance within
the transaction. Build each required persisted index once after loading and
refresh planner statistics. The following independent example assumes that
\texttt{/srv/import/landing\_reads.fasta} contains ACGT sequences with
UniProt-style headers. It uses DNA4 for both the table and the search query.
The accession column deliberately has no ordinary primary-key index.
For a small fixture, place this content in that file on the server:
\par\noindent\begin{minipage}{\linewidth}
\begin{lstlisting}[style=bioseq-appendix-text]
>tr|LDR001|LDR001_HUMAN synthetic read OS=validation_batch PE=1 SV=1
AACCGGTTCTAGTCACTGTAGCCTGATCGTAGCTAGTAACCGGTT
>tr|LDR002|LDR002_MOUSE control read OS=control_batch PE=1 SV=1
TTTTGGGGCCCCAAAATGCATGCATGCATGCATGCATGCATGCA
\end{lstlisting}
\end{minipage}

Load it with:
\par\noindent\begin{minipage}{\linewidth}
\begin{lstlisting}[style=bioseq-appendix-sql]
CREATE TABLE tutorial_loaded_reads (
    accession text NOT NULL,
    organism text NOT NULL,
    seq bioseq.seq('dna4', 64) NOT NULL
) USING bioseq;

BEGIN;
SET LOCAL bioseq.index.insert_maintenance.enabled = off;
\end{lstlisting}
\end{minipage}

\par\noindent\begin{minipage}{\linewidth}
\begin{lstlisting}[style=bioseq-appendix-sql]
SELECT bioseq.load(
    'tutorial_loaded_reads'::regclass,
    '/srv/import/landing_reads.fasta',
    'uni_prot',
    ARRAY[
        ('accession', 'accession')::bioseq.column_mapping,
        ('organism', 'organism_name')::bioseq.column_mapping
    ]
);
\end{lstlisting}
\end{minipage}

\par\noindent\begin{minipage}{\linewidth}
\begin{lstlisting}[style=bioseq-appendix-sql]
SELECT bioseq.update_index(
    'tutorial_loaded_reads', 'minimizers_per_fragment'
);
SELECT bioseq.update_index('tutorial_loaded_reads', 'symbol_counts');
ANALYZE tutorial_loaded_reads;
COMMIT;

SELECT accession, organism, bioseq.len(seq) AS bases
FROM tutorial_loaded_reads
ORDER BY accession;
\end{lstlisting}
\end{minipage}

The loader count is a count of records. The index-update count is a count of
sequence values processed. Neither count alone shows that a later search
uses the intended execution path. Inspect a plan and the returned positions:
\par\noindent\begin{minipage}{\linewidth}
\begin{lstlisting}[style=bioseq-appendix-sql]
BEGIN;
SET LOCAL max_parallel_workers_per_gather = 0;
SET LOCAL bioseq.execution.parallel.enabled = off;
SET LOCAL bioseq.index.sequence_read.enabled = on;

EXPLAIN (ANALYZE, VERBOSE, COSTS OFF)
SELECT accession,
       bioseq.target_start(accession) AS target_start,
       bioseq.target_end(accession) AS target_end
FROM tutorial_loaded_reads
WHERE public.mmsearch(
    seq,
    'CTAGTCACTGTAGCCTGATCGTAGCTAGT'::bioseq.seq('dna4', 64)
)
ORDER BY accession, target_start;
COMMIT;
\end{lstlisting}
\end{minipage}

Disabling both PostgreSQL and BioSeqDB parallel execution makes the
demonstration plan easier to read; it is not required to load a file.
For small ongoing insertions, the
default insertion-time maintenance policy can remain enabled. Transaction
scope restores all settings in these examples automatically.

\subsection{Understand streaming and memory}

The SQL loader sends packed sequence chunks to storage as records are read,
so the complete sequence body need not fit in one PostgreSQL value in memory.
FASTA and FASTQ use one ordered parser state. FASTQ qualities are consumed
to frame records but are not stored in the destination sequence column.

Streaming the packed body does not give a fixed total-memory bound for every
record. Headers, unfinished physical lines, and FASTQ quality scratch space
have separate lifetimes. Quality scratch can grow with the record; a very
long line can reach the reader's buffer limit. The parser should not be used
as a general FASTQ conformance checker.

The setting \texttt{bioseq.load.write\_overlap.enabled} can overlap
parsing with PostgreSQL storage work when the shared scheduler admits a
worker. It is off by default. The optional path preserves record order and
passes owned chunks through a bounded queue. It does not shard a file among
independent SQL parsers. PostgreSQL operations still run on the backend
thread.

\subsection{Distinguish SQL gzip loading from the C++ reader}

Compressed input to \texttt{bioseq.load} uses zlib. The separate C++ batch
reader, \nolinkurl{read_sequences_file}, has additional options: with default
settings it may select a custom parallel deflate decoder for mapped gzip
input of at least 16~MiB. Its default decoded chunk target is 4~MiB; automatic
worker selection uses the process concurrency limit and caps the count at
eight. These are C++ reader policies, not SQL loading settings.

The custom route maps compressed input and returns decompressed chunks in
order. It can fall back to zlib if construction fails before output starts.
A failure after streaming begins is terminal; restarting would risk
duplicating records. Neither route provides gzip range access or a shortcut
to a late record. Aligner, query-cache, and column-compression settings do not
select the gzip backend.

\subsection{Handle transactions and failures}

A load call belongs to its caller's transaction. A parser, mapping, type,
checked constraint, storage, or index-maintenance failure aborts the
statement; it does not commit the accepted prefix. For example, after
creating \texttt{incoming\_reads} above, a fresh batch can be loaded with:
\par\noindent\begin{minipage}{\linewidth}
\begin{lstlisting}[style=bioseq-appendix-sql]
BEGIN;
SET LOCAL bioseq.index.insert_maintenance.enabled = off;

SELECT bioseq.load(
    'incoming_reads'::regclass,
    '/srv/import/run-043.fastq.gz'
);
SELECT bioseq.update_index(
    'incoming_reads', 'minimizers_per_fragment'
);
SELECT bioseq.update_index('incoming_reads', 'symbol_counts');
ANALYZE incoming_reads;
COMMIT;
\end{lstlisting}
\end{minipage}

The server account needs filesystem access, and the database role needs the
applicable database privileges. A client-local path is insufficient when the
server runs elsewhere. Local loading does not automatically retry. After a
failure, roll back the transaction, correct the file or mapping, and start a
new transaction. A successful repeated call appends the file again; the
loader does not deduplicate accessions.

\subsection{Keep local files and remote assemblies separate}

\texttt{bioseq.load} accepts a filesystem path, not an HTTP URL. Remote
relations have their own provider filters, permissions, timeouts, and error
handling. Stage remote rows into a local BioSeqDB table before building local
indexes or repeatedly searching them; see
Section~\ref{sec:appendix-remote-ingestion}. Remote sequence rows currently carry
unbound ASCII values. An ordinary \texttt{INSERT} into a bound local sequence
column validates and encodes them. It uses the ordinary PostgreSQL executor
and the insertion rules supported by the destination table access method.

For the NCBI provider, a single exact \texttt{GCF\_} or \texttt{GCA\_}
assembly accession selects the NCBI Datasets ZIP route. An individual
sequence accession such as \texttt{NC\_045512.2} uses the FASTA retrieval
route. The chunk relation exposes an assembly as bounded output rows:
\par\noindent\begin{minipage}{\linewidth}
\begin{lstlisting}[style=bioseq-appendix-sql]
SELECT chunk_index, base_offset,
       bioseq.len(sequence_chunk) AS bases, end_of_record
FROM bioseq_external.sequence_chunks
WHERE provider = 'ncbi'
  AND accession = 'GCF_009858895.2'
ORDER BY chunk_index;
\end{lstlisting}
\end{minipage}

The ZIP route downloads the complete archive before extraction. Current
limits are 64~MiB of archive bytes, 512~MiB of total decoded content, and
4,096 entries. It rejects unsafe paths, requires the package catalog and
checksum manifest, and selects the catalog's FASTA members. Each selected
member is buffered and MD5-verified before that member emits rows.
Small output rows therefore
avoid one giant SQL value, but do not imply constant memory for an arbitrary
assembly. There is no central-directory range download or late-record seek.
Transient download failures may be retried up to five attempts before
extraction. Once rows have been published, failure terminates the scan.

\subsection{Know the export boundary}

There is no public SQL FASTA, FASTQ, or gzip export function. The C++ writer
can produce plain FASTA and FASTQ, but it is not a SQL export interface and
does not provide gzip output. An application can format results from
\texttt{SELECT} or \texttt{COPY}. To reproduce a source file, it must retain
the required headers, sequence symbols, record order, and, for FASTQ,
qualities separately. Sequence loading alone cannot preserve all of those
details.

The demonstration tables can be removed when no longer needed:
\par\noindent\begin{minipage}{\linewidth}
\begin{lstlisting}[style=bioseq-appendix-sql]
DROP TABLE tutorial_loaded_reads;
DROP TABLE imported_reads;
DROP TABLE incoming_reads;
\end{lstlisting}
\end{minipage}

\section{Generating repeatable sequence tables}
\label{sec:appendix-sequence-generation}
\label{sec:app-sequence-generation}

Two SQL functions create synthetic data directly in \texttt{USING bioseq}
tables. They are useful for query development, repeated examples, and checks
of index and execution paths without downloading data. They return the new
table's OID; casting that value to \texttt{regclass} displays its relation
name. They create tables in the active creation schema, treat the supplied
name as one identifier, and fail on a name collision. They do not append to
or overwrite an existing table. To choose a schema, set \texttt{search\_path}
before calling them rather than embedding a dot in the name.

\subsection{Choose a generator}

\texttt{create\_sample\_table} creates independent, equal-length DNA4
sequences. Its declared columns are \texttt{id BIGSERIAL PRIMARY KEY} and
\texttt{seq bioseq.seq('dna4',32)}. IDs start at one. The alphabet argument
currently accepts only \texttt{dna4}; the width is not configurable.
\par\noindent\begin{minipage}{\linewidth}
\begin{lstlisting}[style=bioseq-appendix-sql]
bioseq.create_sample_table(
    table_name text,
    n_sequences integer,
    sequence_bp_size bigint,
    alphabet text DEFAULT 'dna4',
    seed bigint DEFAULT 42
) RETURNS oid
\end{lstlisting}
\end{minipage}

\texttt{create\_query\_target\_sample\_table} creates a reference with
uniformly sampled A, C, G, and T,
extracts query windows, and applies configured orientation and mutations.
Each row repeats the reference in \texttt{target} and contains one
\texttt{query}. Both sequence columns are declared as
\texttt{bioseq.seq('dna4',64)}. The other columns are \texttt{id},
\texttt{truth\_target\_start}, \texttt{truth\_target\_end},
\texttt{is\_reverse}, \texttt{effective\_error\_rate}, and
\texttt{original\_query\_length}. Its ID column is also a
\texttt{BIGSERIAL PRIMARY KEY} starting at one.
\par\noindent\begin{minipage}{\linewidth}
\begin{lstlisting}[style=bioseq-appendix-sql]
bioseq.create_query_target_sample_table(
    table_name text,
    fragment_count integer,
    fragment_length_mean bigint,
    fragment_length_stddev double precision,
    reference_length bigint,
    substitution_rate double precision,
    insertion_rate double precision,
    deletion_rate double precision,
    insertion_length_mean double precision,
    deletion_length_mean double precision,
    reverse_complement_rate double precision,
    seed bigint
) RETURNS oid
\end{lstlisting}
\end{minipage}

All twelve argument positions of the pair generator must be supplied.
Explicit \texttt{NULL} selects a fallback for any argument except the table
name. In order after the name, those fallbacks are: 100 pairs, mean length
150, length standard deviation 10, reference length 10,000, substitution rate
0.01, insertion and deletion rates 0.0005 each, mean insertion and deletion
lengths 2 each, reverse-complement rate 0.5, and seed 42.

The simple generator accepts 1 to 100,000,000 sequences and a nonnegative
length. Length zero deliberately creates valid empty sequences without a
sequence payload fragment. The pair
generator accepts a positive SQL \texttt{integer} count, a nonnegative mean
fragment length, and a positive reference length. Its count does not share
the simple generator's 100-million check. Use finite, nonnegative
distribution parameters and probabilities in $[0,1]$; the SQL wrapper is not
a comprehensive validator of floating-point arguments. Signed \texttt{bigint}
seeds are accepted.

\subsection{Generate equal-length sequences and inspect symbols}

This call creates 100 independent sequences, each 150 bases long:
\par\noindent\begin{minipage}{\linewidth}
\begin{lstlisting}[style=bioseq-appendix-sql]
SELECT bioseq.create_sample_table(
    'generated_reads', 100, 150, 'dna4', 20260816
)::regclass;

SELECT id, bioseq.len(seq) AS bases,
       bioseq.symbol_count(seq, 'G') AS guanines
FROM generated_reads
ORDER BY id
LIMIT 10;
\end{lstlisting}
\end{minipage}

\texttt{bioseq.len} counts logical symbols independently of the declared
32- or 64-bit storage width. The letters are sampled uniformly from A, C, G,
and T. This table has no
query--target relationship or source-coordinate truth. It is useful for
checking sequence functions and ordinary table workflows.

\subsection{Generate query--target truth pairs}

For reference length $R>1$, the requested Gaussian query length is rounded
and clipped to $[1,R-1]$. For $R=1$, it is one. The source start is sampled
uniformly from positions where that window fits. This is a clipped Gaussian,
not a Gaussian repeatedly sampled until it falls inside the interval.

The generator then makes these choices:
\begin{itemize}
  \item A Bernoulli trial selects whether to reverse-complement the source
        window. Mutation operates on that selected orientation.
  \item Substitution trials apply to every source-fragment symbol.
  \item The indel pass retains each visited symbol, may insert after it,
        and may delete by skipping following source symbols. Skipped symbols
        receive no further indel trials.
  \item Insertion and deletion lengths are Gaussian with the supplied means
        and standard deviation 1.0. They are rounded and clipped at zero,
        so a selected event can have zero length.
\end{itemize}

Truth coordinates are zero-based, half-open positions in the forward target
before mutation, including for reverse queries.
\texttt{original\_query\_length} describes that source interval; the stored
query length can differ after indels. \texttt{effective\_error\_rate} is
exactly the sum of the three configured event rates. It is not a measured
edit fraction, does not account for indel lengths, and can exceed one.

Start with an exact fixture: eight 64-base windows from a 512-base reference,
no substitutions or indels, and a 0.5 probability of reverse orientation.
\par\noindent\begin{minipage}{\linewidth}
\begin{lstlisting}[style=bioseq-appendix-sql]
SELECT bioseq.create_query_target_sample_table(
    'generated_pairs',
    8, 64, 0.0, 512,
    0.0, 0.0, 0.0,
    1.0, 1.0, 0.5, 20260816
)::regclass;

SELECT id, bioseq.len(query) AS query_bases,
       bioseq.len(target) AS target_bases,
       truth_target_start, truth_target_end, is_reverse
FROM generated_pairs
ORDER BY id;
\end{lstlisting}
\end{minipage}

For substitution fixtures, increase \texttt{substitution\_rate} while
leaving both indel rates at zero and inspect alignment scores and edit counts.
For indel fixtures, increase insertion or deletion rates and inspect CIGAR
operations as well as the source interval. A seed alone is insufficient to
describe the experiment: save every generator argument. These generators
make controlled synthetic fixtures, not a general biological simulator.

\subsection{Check repeatability with both generators}

The following example creates two equal-length tables with the same seed,
compares them by ID, and then inspects a small truth table. All three names
must be unused when the example starts.
\par\noindent\begin{minipage}{\linewidth}
\begin{lstlisting}[style=bioseq-appendix-sql]
BEGIN;
SET LOCAL max_parallel_workers_per_gather = 0;
SET LOCAL bioseq.execution.parallel.enabled = off;

SELECT bioseq.create_sample_table(
    'tutorial_generated_a', 4, 48, 'dna4', 314159
)::regclass;
SELECT bioseq.create_sample_table(
    'tutorial_generated_b', 4, 48, 'dna4', 314159
)::regclass;

SELECT a.id, bioseq.len(a.seq) AS bases,
       bioseq.symbol_count(a.seq, 'A') AS adenines,
       a.seq = b.seq AS same_seed_same_sequence
FROM tutorial_generated_a AS a
JOIN tutorial_generated_b AS b USING (id)
ORDER BY a.id;

SELECT bioseq.create_query_target_sample_table(
    'tutorial_generated_pairs',
    3, 48, 0.0, 256,
    0.0, 0.0, 0.0,
    1.0, 1.0, 0.5, 271828
)::regclass;

SELECT id, bioseq.len(query) AS query_bases,
       bioseq.len(target) AS target_bases,
       truth_target_start, truth_target_end, is_reverse,
       effective_error_rate, original_query_length
FROM tutorial_generated_pairs
ORDER BY id;
COMMIT;
\end{lstlisting}
\end{minipage}

The equal-seed comparison should be true for every row. The truth table has
three 48-base queries, 256-base targets, original length 48, and configured
error rate zero. A reverse-complement probability of 0.5 does not require
exactly half of a small sample to be reversed. Repeatability refers to the
same implementation and parameters; Gaussian-derived samples need not be
identical across platforms.

\subsection{Use the normal table lifecycle}

Generated tables support the same subsequent work as other BioSeqDB tables:
\begin{itemize}
  \item Add or update metadata columns for the experiment.
  \item Build persisted indexes with \texttt{bioseq.update\_index}, run
        \texttt{ANALYZE}, and issue \texttt{mmsearch} queries.
  \item Inspect public storage, shard, and index metadata when checking how
        the data was stored and which execution path was selected.
  \item Reuse the same generated inputs for CPU, SIMD, and supported GPU
        search plans, recording the search settings as well as the seed.
\end{itemize}

For example, prepare the independent-sequence table for later searches:
\par\noindent\begin{minipage}{\linewidth}
\begin{lstlisting}[style=bioseq-appendix-sql]
ALTER TABLE generated_reads ADD COLUMN experiment text;
UPDATE generated_reads SET experiment = 'uniform_dna4';
SELECT bioseq.update_index(
    'generated_reads', 'minimizers_per_fragment'
);
SELECT bioseq.update_index('generated_reads', 'symbol_counts');
ANALYZE generated_reads;
\end{lstlisting}
\end{minipage}

The pair generator repeats a reference in every row to keep each truth case
self-contained. For many repeated searches, store targets and queries in
separate tables and use the batch-join workflow instead. This avoids treating
the pair fixture's repeated target payload as the preferred database layout.

\subsection{Understand output, memory, and execution devices}

Both functions write table data inside the server rather than returning a
giant client array. The simple generator streams long sequence payloads into
fragments. Pair generation retains one packed reference and materializes
queries in waves of at most 4,096 before storing them. Memory admission uses
an estimate based on mean query length; it does not cover every Gaussian
tail, insertion expansion, the reference, or all overhead. If reservation
fails, generation can still proceed with one pair. This is not a hard cap on
total memory.

A small SQL result means little client transfer, not little stored data.
Start with small fixtures; the simple generator's maximum row count permits
jobs with substantial storage and runtime requirements. Both SQL generators
use CPU routes. There is no SQL Vulkan random-generation argument, and
aligner Vulkan settings affect later searches rather than these generators.
The generic C++ generator supports more configurations, including alphabets
that these SQL wrappers do not expose.

Synthetic data has no biological provenance. After developing a query,
replace it with local FASTA, FASTQ, or gzip input using
Section~\ref{sec:appendix-local-ingestion}, or with administrator-enabled
external providers using Section~\ref{sec:appendix-remote-ingestion}. Neither SQL
generator is a FASTA or FASTQ export function.

Remove the fixtures after the exercise:
\par\noindent\begin{minipage}{\linewidth}
\begin{lstlisting}[style=bioseq-appendix-sql]
DROP TABLE tutorial_generated_pairs;
DROP TABLE tutorial_generated_b;
DROP TABLE tutorial_generated_a;
DROP TABLE generated_pairs;
DROP TABLE generated_reads;
\end{lstlisting}
\end{minipage}
